\chapter{总结与展望}

本章对全文研究进行收束，概括主要研究发现，并结合数据来源、样本范围、编码方式和平台公开数据限制说明后续研究可以继续改进的方向。

\section{研究总结}

本文以抖音平台省级文旅账号发布的 AI 文旅短视频为研究对象，围绕用户参与行为影响因素展开实证研究。研究以 33 个省份/地区文旅账号中筛选出的 722 条 AI 文旅短视频为分析样本，从技术使用与呈现质量、文旅内容表达、平台观看成本和互动积累时间等层面构建变量体系，并围绕该样本开展描述统计、差异检验、多元回归、稳健性检验、典型案例分析与预测性补充分析。

综合全文结果，可以形成以下几点主要结论。第一，AI 文旅短视频的传播效果具有明显长尾分布特征，少数高互动视频贡献了大部分传播热度。这说明文旅账号即使大量使用 AI 工具，也并不意味着所有内容都能获得稳定高传播，平台中的内容竞争依然十分明显。第二，在控制多个变量后，真正稳定显著影响传播效果的核心因素主要是 AI 技术质感、视频风格、视频时长和发布距采集日天数。其中，AI 技术质感越高，视频越容易获得更高互动；抽象/超现实风格相较纪实/写实风格更容易激发平台用户参与；视频时长越长，互动表现越弱；发布距采集日天数越长，累计互动水平越高。第三，AI 介入形式、AI 介入程度、视频主体、视频主题和内容导向在单因素层面虽然存在差异，但在控制其他因素后并未表现出稳定独立解释力，说明这些变量更多反映内容类型差异，而非最核心的传播驱动因素。

进一步结合典型案例和预测性补充分析，可以看出本文结论具有较强一致性。高互动样本往往同时具备较高技术质感、较鲜明风格和较紧凑节奏，而低互动样本则更容易出现画面粗糙、表达平直或平台感不足的问题。机器学习模型中的变量重要性排序与回归分析高度相近，进一步说明本文识别出的核心因素并非某一种分析方法下的偶然结果，而是在解释与预测两个层面均具有较强稳定性。

\section{研究展望}

本文仍存在一定研究边界，主要来自平台公开数据本身的限制。本文使用的是抖音平台公开可见数据，无法获得曝光量、完播率、点击率、平均观看时长、粉丝画像、推荐链路和投流情况等后台指标，因此只能以点赞、评论、收藏和转发等公开互动数据衡量传播效果。这些指标能够反映用户外显互动，但不能完整呈现平台推荐、观看留存和互动转化的全过程。

本文基于单次采集截面的累计互动数据展开分析。虽然研究中已经统一采集日期，并在模型中控制发布距采集日天数，也通过日均互动量进行稳健性检验，但短视频互动仍然具有持续累积和阶段性扩散特征。单次截面数据更适合分析内容特征与累计互动水平之间的关系，难以完整呈现视频发布后不同阶段的传播变化。

在人工编码方面，本文已经通过统一编码规则、结果变量隔离、样本复核和保守归类等方式降低主观判断误差，但 AI 技术质感、视频风格和内容导向等变量仍带有一定内容判断属性。本文尝试使用 AI 对部分样本进行辅助复核，以识别高分歧变量和修订编码边界，但未将其作为正式双人一致性检验结果报告。后续研究若具备条件，仍应优先采用双人独立编码与一致性统计，以进一步提高内容分析结果的可靠性。

本文使用的是观察性平台数据，平台推荐机制、热点事件、账号粉丝基础、发布时间段和投流情况等因素无法被完整观测。因此，本文结论更适合作为内容特征与传播效果之间的相关关系和影响因素识别，而不是严格意义上的因果效应判断。后续如果能够获得更完整的后台数据或连续追踪数据，可以进一步区分曝光、观看、互动和转化等不同传播环节。

总体来看，本文已经在公开数据条件下尽量统一样本口径、控制时间因素并进行稳健性检验，但 AI 文旅短视频传播效果仍受到平台机制和数据可得性的客观限制。未来研究若能在数据权限和连续观测条件上有所改善，将有助于更细致地解释 AI 内容从平台推荐到用户互动的具体过程。
